\section{Markowitz como caso base}

El caso base del proyecto es el portafolio \textbf{Markowitz} de optimizaci\'on
media-varianza. Todas las metodolog\'ias basadas en Black-Litterman se comparan
contra esta referencia, de modo que una \textit{view} solo se considera valiosa
si mejora el desempe\~no del cliente respecto de Markowitz bajo el mismo
protocolo de simulaci\'on.

La elecci\'on de Markowitz como caso base se justifica porque es la l\'inea base
anal\'itica natural del problema: optimiza expl\'icitamente la relaci\'on
riesgo-retorno a partir de las estimaciones hist\'oricas de $\mu$ y $\Sigma$,
sin incorporar \textit{views} adicionales. As\'i, aislar la contribuci\'on de
Black-Litterman equivale a medir cu\'anto aporta el prior de equilibrio y las
\textit{views} calibradas por sobre la estimaci\'on hist\'orica pura. En
versiones anteriores del proyecto se utiliz\'o un benchmark equiponderado como
caso base; este se descart\'o por ser una regla de asignaci\'on trivial que no
representa una alternativa de inversi\'on optimizada y comparable.

\section{M\'etricas del caso base}

La Tabla~\ref{tab:markowitz_base_kpis} resume el desempe\~no de Markowitz base
en el horizonte de test limpio P4, agregado sobre perfiles, bajo el modelo de
abandono semanal y la comisi\'on de 0,5\% mensual sobre el saldo activo. Al ser
independiente de \textit{views}, Markowitz base entrega las mismas m\'etricas en
todas las comparaciones por \textit{view}.

\begin{table}[h]
\centering
\caption{Desempe\~no de Markowitz base en test limpio P4 (agregado sobre
perfiles, comisi\'on 0,5\% mensual sobre saldo).}
\label{tab:markowitz_base_kpis}
\small
\begin{tabular}{@{}lrr@{}}
\toprule
Indicador & Sin pandemia & Con pandemia \\
\midrule
Riqueza terminal media (USD) & 2\,342 & 2\,140 \\
Clientes retenidos (de 1\,000) & 865 & 850 \\
Turnover medio & 9,0\% & 9,4\% \\
HHI sectorial & 0,182 & 0,178 \\
Utilidad FinPUC por saldo (USD) & 258 & 243 \\
\bottomrule
\end{tabular}
\end{table}

Estos valores constituyen la referencia frente a la cual se mide la mejora de
cada \textit{view} de Black-Litterman. Una metodolog\'ia agrega valor cuando
incrementa la riqueza terminal del cliente y la retenci\'on sin elevar de forma
excesiva el turnover ni la concentraci\'on sectorial; el efecto neto sobre la
utilidad de FinPUC se sigue, como consecuencia, del mayor saldo activo
administrado que esa mejora sostiene.
